\documentclass[11pt,a4paper]{article}
% preamble.tex — estilo común de todos los apuntes Froth.
% Cada apunte hace % preamble.tex — estilo común de todos los apuntes Froth.
% Cada apunte hace % preamble.tex — estilo común de todos los apuntes Froth.
% Cada apunte hace \input{preamble} justo después de \documentclass.
\usepackage[margin=2.4cm]{geometry}
\usepackage{xcolor}
\definecolor{litblue}{RGB}{31,79,121}
\definecolor{litgreen}{RGB}{27,94,32}
\definecolor{litorange}{RGB}{230,81,0}
\usepackage{parskip}
\usepackage{titlesec}
\titleformat{\section}{\Large\bfseries\color{litblue}}{\thesection}{0.6em}{}
\titleformat{\subsection}{\large\bfseries\color{litblue}}{\thesubsection}{0.6em}{}
\usepackage{tcolorbox}
\usepackage{fancyhdr}
\pagestyle{fancy}
\fancyhf{}
\fancyhead[L]{\small\color{litblue}Froth — Apuntes de ML}
\fancyhead[R]{\small\thepage}
\renewcommand{\headrulewidth}{0.4pt}
\usepackage[colorlinks=true,linkcolor=litblue,urlcolor=litblue]{hyperref}

% Cabecera de cada apunte: \cabecera{numero}{titulo}
\newcommand{\cabecera}[2]{%
  \begin{tcolorbox}[colback=litblue,coltext=white,colframe=litblue,arc=2mm]
    {\Large\bfseries Apunte #1 — #2}\\[3pt]
    {\small Proyecto Froth · aprender Machine Learning construyendo}
  \end{tcolorbox}\vspace{0.8em}}

% Cajas temáticas
\newtcolorbox{concepto}[1]{colback=blue!4,colframe=litblue,title={Concepto: #1},arc=1.5mm}
\newtcolorbox{practica}{colback=green!5,colframe=litgreen,title={Pruébalo tú (teclado en mano)},arc=1.5mm}
\newtcolorbox{ojo}{colback=orange!8,colframe=litorange,title={Ojo / error típico},arc=1.5mm}
\newtcolorbox{resumen}{colback=gray!8,colframe=gray!60,title={En una frase},arc=1.5mm}

\newcommand{\codigo}[1]{\texttt{#1}}
 justo después de \documentclass.
\usepackage[margin=2.4cm]{geometry}
\usepackage{xcolor}
\definecolor{litblue}{RGB}{31,79,121}
\definecolor{litgreen}{RGB}{27,94,32}
\definecolor{litorange}{RGB}{230,81,0}
\usepackage{parskip}
\usepackage{titlesec}
\titleformat{\section}{\Large\bfseries\color{litblue}}{\thesection}{0.6em}{}
\titleformat{\subsection}{\large\bfseries\color{litblue}}{\thesubsection}{0.6em}{}
\usepackage{tcolorbox}
\usepackage{fancyhdr}
\pagestyle{fancy}
\fancyhf{}
\fancyhead[L]{\small\color{litblue}Froth — Apuntes de ML}
\fancyhead[R]{\small\thepage}
\renewcommand{\headrulewidth}{0.4pt}
\usepackage[colorlinks=true,linkcolor=litblue,urlcolor=litblue]{hyperref}

% Cabecera de cada apunte: \cabecera{numero}{titulo}
\newcommand{\cabecera}[2]{%
  \begin{tcolorbox}[colback=litblue,coltext=white,colframe=litblue,arc=2mm]
    {\Large\bfseries Apunte #1 — #2}\\[3pt]
    {\small Proyecto Froth · aprender Machine Learning construyendo}
  \end{tcolorbox}\vspace{0.8em}}

% Cajas temáticas
\newtcolorbox{concepto}[1]{colback=blue!4,colframe=litblue,title={Concepto: #1},arc=1.5mm}
\newtcolorbox{practica}{colback=green!5,colframe=litgreen,title={Pruébalo tú (teclado en mano)},arc=1.5mm}
\newtcolorbox{ojo}{colback=orange!8,colframe=litorange,title={Ojo / error típico},arc=1.5mm}
\newtcolorbox{resumen}{colback=gray!8,colframe=gray!60,title={En una frase},arc=1.5mm}

\newcommand{\codigo}[1]{\texttt{#1}}
 justo después de \documentclass.
\usepackage[margin=2.4cm]{geometry}
\usepackage{xcolor}
\definecolor{litblue}{RGB}{31,79,121}
\definecolor{litgreen}{RGB}{27,94,32}
\definecolor{litorange}{RGB}{230,81,0}
\usepackage{parskip}
\usepackage{titlesec}
\titleformat{\section}{\Large\bfseries\color{litblue}}{\thesection}{0.6em}{}
\titleformat{\subsection}{\large\bfseries\color{litblue}}{\thesubsection}{0.6em}{}
\usepackage{tcolorbox}
\usepackage{fancyhdr}
\pagestyle{fancy}
\fancyhf{}
\fancyhead[L]{\small\color{litblue}Froth — Apuntes de ML}
\fancyhead[R]{\small\thepage}
\renewcommand{\headrulewidth}{0.4pt}
\usepackage[colorlinks=true,linkcolor=litblue,urlcolor=litblue]{hyperref}

% Cabecera de cada apunte: \cabecera{numero}{titulo}
\newcommand{\cabecera}[2]{%
  \begin{tcolorbox}[colback=litblue,coltext=white,colframe=litblue,arc=2mm]
    {\Large\bfseries Apunte #1 — #2}\\[3pt]
    {\small Proyecto Froth · aprender Machine Learning construyendo}
  \end{tcolorbox}\vspace{0.8em}}

% Cajas temáticas
\newtcolorbox{concepto}[1]{colback=blue!4,colframe=litblue,title={Concepto: #1},arc=1.5mm}
\newtcolorbox{practica}{colback=green!5,colframe=litgreen,title={Pruébalo tú (teclado en mano)},arc=1.5mm}
\newtcolorbox{ojo}{colback=orange!8,colframe=litorange,title={Ojo / error típico},arc=1.5mm}
\newtcolorbox{resumen}{colback=gray!8,colframe=gray!60,title={En una frase},arc=1.5mm}

\newcommand{\codigo}[1]{\texttt{#1}}

\begin{document}
\cabecera{43}{La puerta de relevancia: separar el grano de la paja en 190.000 papers}

\section{Glosario en palabras simples (léelo primero)}

\begin{itemize}
  \item \textbf{Vector}: la ``huella digital'' numérica de un paper --- una lista de 768
        números que captura su significado. Papers que hablan de lo mismo tienen huellas
        parecidas (apunte 05).
  \item \textbf{Coseno}: el medidor de parecido entre dos huellas. Va de $-1$ a $1$:
        cerca de $1$ = hablan de lo mismo; cerca de $0$ = nada que ver.
  \item \textbf{Centroide}: el ``promedio'' de un grupo de vectores --- su centro de
        gravedad, como la ley media de un lote de mineral.
  \item \textbf{Umbral}: el valor de corte. Encima: el paper entra al mapa. Debajo: fuera.
  \item \textbf{Bimodal}: distribución con DOS jorobas. Si los buenos hacen una y los
        malos otra, el valle entre ambas es el corte natural.
  \item \textbf{Hub / cluster}: grupo de papers muy parecidos entre sí --- un subtema.
\end{itemize}

\section{El problema que TÚ encontraste}

Con la cosecha sin techo (190.000 papers), abriste el mapa de tailings y viste
\emph{``cosas hasta de recursos humanos, en colores''}. Tu hipótesis fue exacta: ninguna
palabra falló al vectorizarse --- falló \textbf{la puerta} que decide qué entra al mapa.

La puerta vieja: relevancia $= \cos(\mbox{paper},\ \mbox{TÍTULO del topic})$, con
umbral fijo 0.55. Medido en tailings: aceptaba el \textbf{93\%} de 16.681 papers,
incluyendo clamidiosis en koalas (por ``iron'' biológico), ``event management'' (por
``management'') y ``Circular Economy in SMEs'' (por ``circular economy''). Una frase
corta con imanes genéricos no puede distinguir facetas. Además, el techo viejo de 25
resultados por término era un filtro de precisión ACCIDENTAL: al quitarlo (decisión
correcta: cobertura es la misión), entró la cola profunda de cada término.

\section{Mini-ejemplo de juguete (haz las cuentas conmigo)}

Huellas de solo 2 números (en vez de 768). El título apunta a $T=(1,\ 0)$. Tres papers:

\begin{center}
\begin{tabular}{lcc}
 & vector & $\cos$ con el título \\
\hline
A: flotación de relaves & $(0.98,\ 0.20)$ & $0.98$ \\
B: gestión de personal  & $(0.60,\ 0.80)$ & $0.60$ \\
C: koalas               & $(0.50,\ 0.87)$ & $0.50$ \\
\end{tabular}
\end{center}

Con umbral 0.55, ¡B entra al mapa! ($0.60 > 0.55$). Ese es tu hub de recursos humanos.

La \textbf{puerta contrastiva} pregunta algo mejor: ¿te pareces a mi tema MÁS de lo que
te pareces a lo genérico?
\[
v_2 = \cos(\mbox{paper},\ \mbox{núcleo}) - 0.4 \cdot \cos(\mbox{paper},\ \mbox{global})
\]
donde núcleo $\approx A$ (los papers más pegados al título) y global $=$ promedio de
todos $\approx (0.74,\ 0.67)$ normalizado. Cuentas:
\[
v_2(A) = 1.00 - 0.4 \cdot 0.86 = \mathbf{0.66}
\qquad
v_2(B) = 0.75 - 0.4 \cdot 0.98 = \mathbf{0.36}
\]
B queda lejos de cualquier corte razonable. Con 768 números y 16.681 papers, esto mismo
produjo una distribución \textbf{bimodal} con valle en $+0.399$ --- y el umbral se elige
en ESE valle, por corpus (perilla calculada, no constante fósil).

\begin{concepto}{la idea contrastiva}
Es el mismo truco que arregló los títulos KeyBERT (apunte 40): premiar cercanía al
núcleo Y castigar cercanía a la masa genérica. Lo genérico compartido deja de contar
como mérito.
\end{concepto}

\section{Tropiezo 1: la puerta por-paper aplanó los mapas}

La primera versión juzgaba paper por paper\ldots{} y tailings pasó de 115 subtemas a
\textbf{2} (a CUALQUIER granularidad --- lo confirmó el barrido DBCV). ¿Por qué? Al
quedarte solo con ``los parecidos al núcleo'', la población superviviente es una banda
homogénea: una bola sin valles internos que HDBSCAN no puede subdividir. Ganamos
precisión, destruimos la ESTRUCTURA --- que es la mitad del producto.

El arreglo salió de tu propia frase (``un hub de RRHH, \emph{en colores}''): la basura
llega en \textbf{hubs coherentes}, así que la puerta debe juzgar HUBS: (1) mapear TODO,
(2) puntuar cada cluster por el $v_2$ promedio de sus miembros, (3) botar clusters
enteros bajo el valle de esas medias, (4) el ruido suelto se juzga individual. Medido:
los hubs de física de partículas, astronomía, RRHH y koalas murieron completos; el de
morteros con relaves no perdió ni un paper; 325 subtemas del dominio intactos.

\section{Tropiezo 2: tu doble ciego destapó un sesgo}

Juzgaste 40 casos de frontera a ciegas: 62\% de acuerdo, pero con dirección --- rescataste
11 papers que la máquina botaba, casi todos de \textbf{economía circular}\ldots{} que está
EN el título de tu tesis. Causa: el centro global apunta en parte hacia el propio tema
(el corpus está lleno de él), así que un paper alineado con el título recibía premio del
núcleo Y castigo del global: doble penalización por su propia identidad. Arreglo:
\textbf{proyectar fuera} del global la dirección del núcleo antes de castigar. Verificado:
``Towards Circular Economy Metrics'' volvió; koalas y el paper del queso (``Iron
Ages''\ldots{} de la Edad del Hierro, no del mineral) siguen fuera.

Honestidad de cierre: re-medido tu acuerdo contra la puerta final, quedó en 60\% ---
prácticamente igual. El fix ganó los casos que te importaban (los de economía circular),
pero la frontera de $\pm0.04$ del valle es ambigua HASTA para el experto (tú dejaste
pasar el queso por el ``Iron''). No se afina más: sería sobreajustar a 40 casos.

\begin{ojo}
El valor del gate está en el nivel de HUBS (basura coherente muere entera, estructura
del dominio intacta), no en decorar el borde. Ninguna frontera unidimensional es
perfecta; el flujo correcto es máquina propone $\to$ experto valida $\to$ la perilla se
recalibra donde el experto tiene razón sistemática (economía circular sí, queso no).
\end{ojo}

\begin{resumen}
Puerta vieja: $\cos$ con el título + umbral fijo = 93\% aceptado, koalas en colores.
Puerta nueva: score contrastivo (núcleo vs.\ global SIN la dirección del tema) $\to$
bimodal $\to$ valle POR corpus $\to$ juicio por HUBS para preservar estructura.
Resultado: 22/22 topics limpios con su propio umbral, master de 77.636 papers únicos
relevantes. Tres lecciones: (1) contrastar contra lo genérico separa lo que el parecido
a secas no puede; (2) filtrar individuos puede destruir la estructura colectiva ---
juzga al grupo cuando la señal es grupal; (3) el juez humano encuentra los sesgos que
la métrica no ve --- y también tiene los suyos.
\end{resumen}

\end{document}
